\documentclass{article}%
\usepackage[T1]{fontenc}%
\usepackage[utf8]{inputenc}%
\usepackage{lmodern}%
\usepackage{textcomp}%
\usepackage{lastpage}%
\usepackage{geometry}%
\geometry{margin=2.2cm}%
%
\usepackage{amsmath}%
\usepackage{amssymb}%
\usepackage{booktabs}%
\usepackage{longtable}%
\usepackage{array}%
\usepackage{hyperref}%
\hypersetup{hidelinks}%
\title{ITA 2026 {-} Química (claude\_haiku\_5\_5) x Poliedro\\[4pt]\large Comparação com o gabarito Poliedro Resolve\thanks{Gerado com gpt-resolve: \url{https://github.com/lgabs/gpt-resolve}}}%
\date{10/11/2026}%
%
\begin{document}%
\normalsize%
\maketitle%
\section*{Resumo}\textbf{Nota média: 10,0/10} (32 de 32 itens com a resolução correta; cada item vale o mesmo dentro da questão).\par\medskip%
\begin{longtable}{@{}c c p{0.78\textwidth}@{}}\toprule Questão & Nota & Resumo \\ \midrule
1 & 10,0 & Questão 1: a resolução e o gabarito do Poliedro coincidem nos três itens: a) 0,4 mol/L; b) 2,0 mol/L; c) 0,8/38 $\approx$ 2,1$\times$10$^{-}$$^2$ mol/L. Refiz as contas em Python e, com c2 = 0,4 mol/L, o balanço de cargas completo devolve pH = 12,000. A resolução usa a condição protônica, o gabarito usa a tabela de hidrólise, e os dois caminhos são equivalentes e corretos. A resolução evita o erro de usar a aproximação $\surd$(Kh$\cdot$C), que daria 0,2 mol/L, e confere tudo no final. Não há divergências. Restam só pontos menores: Kw não está na tabela de constantes da prova, a 'equivalência' tem sentido apenas estequiométrico com pKA = 12, e há erros tipográficos no enunciado e no gabarito. \\
2 & 10,0 & A resolução está correta e concorda com o gabarito do Poliedro nos dois itens; as diferenças são só de arredondamento. No item a, os dois obtêm n(O2) = 0,025 mol e V $\approx$ 0,61 L: a resolução chega a 0,612 L com o R da prova (0,0821) e o gabarito a 0,6109 L com R = 0,082. No item b, os dois usam primeira ordem, e o valor exato é 0,5$\cdot$e\^{}(-0,3465) $\approx$ 0,3536 mol/L. A resolução arredonda para 0,354 e o gabarito para 0,353 (0,707$^3$). O gabarito tem dois deslizes menores: escreve a unidade 'mol' em vez de mol/L e usa R = 0,082 em vez do valor fornecido. \\
3 & 10,0 & A resolução está correta nos dois itens e concorda com o gabarito Poliedro, a menos de arredondamentos. a) A resolução dá 552,8 s (exato: 1000(1-1/$\surd$5)) e o gabarito 553,6 s; a diferença vem só de o gabarito usar $\surd$5 = 2,24. b) A resolução dá {[}A?{]} $\approx$ 6,11$\cdot$10$^{-}$$^3$ e {[}B$^{-}${]} $\approx$ 1,64$\cdot$10$^{-}$$^3$ mol/L, valores exatos que fecham K\_PS = 10$^{-}$$^{5}$. O gabarito dá 6,135$\cdot$10$^{-}$$^3$ e 1,655$\cdot$10$^{-}$$^3$ mol/L porque arredondou o termo constante da equação do 2º grau (3,014$\cdot$10$^{-}$$^{5}$ $\to$ 3$\cdot$10$^{-}$$^{5}$), e esses valores dão produto 1,014$\cdot$10$^{-}$$^{5}$. Nesse ponto a resolução é mais precisa que o gabarito, mas o método e a conclusão são os mesmos. \\
4 & 10,0 & A resolução está correta nos dois itens e concorda com o Poliedro Resolve. No item a), ela dá 5,58 g (cerca de 5,6 g) e o gabarito dá 5,57 g. A diferença de 0,2\% vem só de usar o resíduo medido de 28,5\% em vez do valor calculado com n = 9 (28,57\%), e a própria resolução mostra os dois resultados (5,58 g e 5,56 g). No item b), os dois chegam a n = 9 (calculado n = 9,06). O único problema encontrado está no gabarito: a equação 0,285 = 266/(303 + 18n) está sem o fator 2 no denominador. É erro de digitação e não muda a resposta. \\
5 & 10,0 & A resolução está correta nos quatro itens e concorda com o gabarito do Poliedro. a) São 3 arranjos, um para cada átomo central (C, N ou S). b) BF3 e NH3, ambas entre as respostas aceitas pelo gabarito. c) Os produtos são o orto e o para{-}fluoronitrobenzeno, com o para majoritário. d) A ordem é H2S < CH3OH < H2O. Conferi o enunciado no PDF original (página 7) e refiz as cargas formais, os balanceamentos e as massas molares: não há erros. Os únicos pontos em aberto são de interpretação ou de dados citados de memória, nenhum deles de peso. \\
6 & 10,0 & A resolução está correta e coincide com o gabarito do Poliedro nos três itens: a) não favorece, b) não favorece, c) favorece. Os pK derivados também são os mesmos: pKa(H3O+) = -1,70, pKb(propinido) = -11, pKa(NH4+) = 9,25 e pKa(NH3) = 36. Conferi as constantes de equilíbrio em Python: 10\^{}{-}26,7 $\approx$ 2,0$\times$10\^{}{-}27, 10\^{}{-}15,75 $\approx$ 1,8$\times$10\^{}{-}16 e 10\^{}11. Na resolução, o propinido aparece como CH3C$\equiv$C- e o propino como CH3C$\equiv$CH. A transcrição do enunciado confere com a página 8 da prova. Ficam só observações de convenção: o pKb da água dado e os valores que a resolução acrescentou por hipótese para H2O/OH-. Nenhuma delas afeta as conclusões. \\
7 & 10,0 & A resolução e o gabarito do Poliedro coincidem em método e resultado nos cinco itens. As pequenas diferenças numéricas vêm de o gabarito usar N\_A = 6$\times$10$^2$$^3$ e M(H$_{2}$) = 2 g/mol, enquanto a resolução usa os dados da prova (6,02$\times$10$^2$$^3$ e 2,02 g/mol): 5,32$\times$10$^{5}$ J contra 5,304$\times$10$^{5}$ J, 1,35$\times$10$^{-}$$^2$ contra 1,36$\times$10$^{-}$$^2$ mol, 1,49$\times$10$^{9}$ J contra 1,50$\times$10$^{9}$ J, 837 contra 835 nm. Refiz as contas em Python e os valores da resolução estão corretos com as constantes fornecidas. No item (a), as duas trazem semirreações válidas. O gabarito mistura meio básico (cátodo) e ácido (ânodo), mas a soma ainda fecha. A questão está correta: 5/5 itens equivalentes ao gabarito. \\
8 & 10,0 & A resolução está correta e bate com o Poliedro Resolve nos quatro itens. a) CH4 + 2 O2 $\to$ CO2 + 2 H2O(g), $\Delta$H = -800 kJ/mol. b) 960 g de O2 e 240,75 g de CH4 contra 240 g no gabarito; a diferença vem só de M(CH4) = 16,05 contra 16. c) 0,644 g contra 0,64 g, mesma diferença de arredondamento. d) 200 g nos dois. Confirmei todas as contas em Python e conferi o enunciado no PDF original. Nenhum dos dois tem erro. As únicas diferenças são de arredondamento e de apresentação (o gabarito também escreve a equação multiplicada por 15, com $\Delta$H = -12 MJ). \\
9 & 10,0 & Questão 9: a resolução concorda com o gabarito Poliedro nos quatro itens. No item a ela lista as mesmas espécies e acrescenta o HCrO$_{4}$$^{-}$; a única diferença é que põe o HSO$_{4}$$^{-}$ entre as espécies principais, enquanto o gabarito o considera minoritário. Nesse ponto a resolução está quimicamente mais correta, pois Ka2 $\approx$ 10$^{-}$$^2$ torna o HSO$_{4}$$^{-}$ majoritário em H$_{2}$SO$_{4}$ diluído usual. Nos demais itens o conteúdo é o mesmo: b) as camadas se formam por migração iônica; c) Pt$_{1}$ é o cátodo (2H? + 2e$^{-}$ $\to$ H$_{2}$) e Pt$_{2}$ é o ânodo (2H$_{2}$O $\to$ O$_{2}$ + 4H? + 4e$^{-}$); d) P$_{1}$ é negativo, P$_{2}$ é positivo e a corrente convencional vai de P$_{2}$ para P$_{1}$ pelo circuito externo. As contas extras conferem: $\Delta$G$^\circ$ $\approx$ +4,75$\times$10$^{5}$ J/mol e fração de HCrO$_{4}$$^{-}$ de 50\%. Não há erro na resolução. \\
10 & 10,0 & A resolução está certa e coincide com o gabarito do Poliedro Resolve nos três itens: a) X = 226, b) Z = 86, c) 3 partículas alfa. Conferi a transcrição com o enunciado original (página 12) e refiz a cadeia de decaimentos: Ac{-}232 emite $\beta$$^{-}$ e vira Th{-}232 (142 nêutrons), E é Po{-}226 e A é Rn{-}226, formado por 3 emissões alfa a partir do U{-}238. O raciocínio é o mesmo do gabarito, com as hipóteses de leitura declaradas. Não há divergência. \\
\bottomrule\end{longtable}%
\section*{Principais pontos de dúvida}%
\begin{itemize}%
\item%
\textbf{Q4} \textit{(key)} No item b, o gabarito escreve 0,285 = 266/(303 + 18n), sem o fator 2 no denominador. O correto é 0,285 = 266/{[}2(303 + 18n){]}. Lida literalmente, a equação dá n $\approx$ 35. O texto do gabarito ('2 mols de oxalato $\to$ 1 mol de Nb2O5') e a resposta n = 9 mostram que é erro de digitação, mas a equação como está escrita não leva ao resultado.%
\end{itemize}%
\section*{Questão 1}%
\href{https://poliedroresolve.sistemapoliedro.com.br/vestibulares/ita/2026/ita-2-fase-dia-2-29-10-2025-2-fase-dia-2/dissertativa-1-quimica-2-fase-dia-2-ita-2026}{Resolução do gabarito}\par\smallskip%
\begin{longtable}{@{}p{0.06\textwidth} p{0.33\textwidth} p{0.33\textwidth} p{0.1\textwidth} p{0.08\textwidth}@{}}\toprule Item & Resolução & Gabarito & Veredito & Certo \\ \midrule
a & c2 = 0,40 mol/L & (M\_A)2 = 0,4 mol/L & igual & ambos \\
b & c1 = 2,0 mol/L & (M\_A)1 = 2,0 mol/L & igual & ambos \\
c & c3 = 2/95 $\approx$ 2,1$\times$10$^{-}$$^2$ mol/L & M\_B = 0,8/38 $\approx$ 2,1$\cdot$10$^{-}$$^2$ mol/L & igual & ambos \\
\bottomrule\end{longtable}%
\begin{itemize}%
\item%
\textbf{a)} Os dois fazem a mesma conta: C = c2$\cdot$2/40 = c2/20 = 5$\cdot$10$^{-}$$^2$$\cdot$c2 na mistura. Com pH = 12, {[}OH$^{-}${]} = {[}HA{]} = 10$^{-}$$^2$, e Kh = Kw/KA = 10$^{-}$$^2$. Então Kh = 10$^{-}$$^{4}$/(C - 10$^{-}$$^2$) = 10$^{-}$$^2$ dá C = 2$\cdot$10$^{-}$$^2$, logo c2 = 0,4 mol/L. A resolução chega ao mesmo resultado pela condição protônica junto com {[}A$^{-}${]}/{[}HA{]} = KA/{[}H?{]} = 1, já que pH = pKA. Conferi numericamente: com c2 = 0,4, o balanço de cargas completo dá pH = 12,000.%
\item%
\textbf{b)} Diluição de 4 mL para 20 mL, fator 5: c1 = 5$\cdot$0,4 = 2,0 mol/L.%
\item%
\textbf{c)} Estequiometria 1:1: 0,4$\cdot$2 = c3$\cdot$38, logo c3 = 0,8/38 = 2/95 $\approx$ 0,02105 mol/L. A fração 2/95 é igual a 0,8/38, e as duas formas decimais coincidem.%
\end{itemize}%
\paragraph{Pontos de dúvida}%
\begin{itemize}%
\item%
\textit{(solution)} Kw = 1,0$\times$10$^{-}$$^{1}$$^{4}$ a 25 $^\circ$C não aparece na tabela de constantes da prova. A resolução e o gabarito usam esse valor sem justificar. É a hipótese padrão e a resolução a declara.%
\item%
\textit{(statement)} Como pH = pKA = 12, no 'ponto de equivalência' metade do ânion está hidrolisado a HA. A equivalência aqui é só estequiométrica (n ácido = n base), porque na prática uma titulação com pKA = 12 não teria ponto final nítido. Isso não muda a conta, mas é uma sutileza conceitual do enunciado.%
\item%
\textit{(statement)} O enunciado traz o erro de digitação 'Solução I' no lugar de 'Solução 1'. A leitura é inequívoca, e os dois interpretam igual.%
\item%
\textit{(solution)} Armadilha comum: usar a fórmula simplificada {[}OH$^{-}${]} $\approx$ $\surd$(Kh$\cdot$C) leva a c2 = 0,20 mol/L, metade do valor correto, porque a hidrólise de 50\% não é desprezível. A resolução aponta isso de forma explícita. Pode interessar ao professor como critério de correção.%
\item%
\textit{(key)} A resolução do gabarito tem um defeito de formatação na equação ('HA + BOH $\to$ \^{}BA + H$_{2}$O'). É só tipografia e não afeta o resultado.%
\end{itemize}%
\section*{Questão 2}%
\href{https://poliedroresolve.sistemapoliedro.com.br/vestibulares/ita/2026/ita-2-fase-dia-2-29-10-2025-2-fase-dia-2/dissertativa-2-quimica-2-fase-dia-2-ita-2026}{Resolução do gabarito}\par\smallskip%
\begin{longtable}{@{}p{0.06\textwidth} p{0.33\textwidth} p{0.33\textwidth} p{0.1\textwidth} p{0.08\textwidth}@{}}\toprule Item & Resolução & Gabarito & Veredito & Certo \\ \midrule
a & V(O2) $\approx$ 0,612 L ($\approx$ 612 mL), com n(O2) = 0,025 mol, R = 0,0821 atm$\cdot$L/(K$\cdot$mol), T = 298 K, P = 1 atm & V(O2) $\approx$ 0,6109 L (n(O2) = 0,025 mol; PV = nRT com R = 0,082, a 298 K e 1 atm) & equivalente & ambos \\
b & {[}H2O2{]} $\approx$ 0,354 mol/L (0,5$\cdot$e\^{}(-0,3465) = 0,5$\cdot$0,707) & {[}H2O2{]} $\approx$ 0,353 mol (deveria ser mol/L), calculado como (e\^{}(-0,3465))\^{}3 = 0,707\^{}3 & equivalente & ambos \\
\bottomrule\end{longtable}%
\begin{itemize}%
\item%
\textbf{a)} O raciocínio é o mesmo: n(H2O2) = 0,5$\cdot$0,1 = 0,050 mol, o que dá n(O2) = 0,025 mol, e V = nRT/P. A diferença está só no valor de R. A resolução usou R = 0,0821, que é o valor dado na prova: 0,025$\cdot$0,0821$\cdot$298 = 0,6116 L. O gabarito usou R = 0,082 e obteve 0,6109 L. Com R = 8,31 J/(K$\cdot$mol) sai 0,6110 L. As três contas diferem cerca de 0,1 \%, que é só arredondamento.%
\item%
\textbf{b)} Os dois usam primeira ordem (k em min$^{-}$$^{1}$) e {[}A{]} = {[}A{]}0$\cdot$e\^{}(-kt) com kt = 0,3465. O valor exato é 0,5$\cdot$e\^{}(-0,3465) = 0,35358 = 2\^{}(-3/2). O caminho do gabarito também é válido, porque 0,5 = e\^{}(-0,693) = (e\^{}(-0,3465))\^{}2, e então 0,5$\cdot$e\^{}(-0,3465) = (e\^{}(-0,3465))\^{}3. A diferença está só no arredondamento: 0,5$\cdot$0,707 = 0,3535, que vira 0,354, e 0,707\^{}3 = 0,3534, que vira 0,353. O gabarito escreveu a unidade como 'mol' em vez de mol/L. É um deslize de notação.%
\end{itemize}%
\paragraph{Pontos de dúvida}%
\begin{itemize}%
\item%
\textit{(key)} No item b o gabarito dá a resposta em 'mol', mas o enunciado pede concentração, que é em mol/L. É um deslize de unidade e não muda o valor numérico.%
\item%
\textit{(key)} No item a o gabarito usa R = 0,082 atm$\cdot$L/(K$\cdot$mol), mas a prova fornece R = 8,21$\times$10$^{-}$$^2$ (0,0821). Por isso o gabarito chega a 0,6109 L e a resolução a 0,6116 L, uma diferença de 0,1 \% que é só arredondamento.%
\item%
\textit{(key)} No item b a cadeia de igualdades do gabarito está escrita de forma confusa: aparece '= -0,693 = -10$\cdot$0,0693' no meio de uma expressão que vale -1,0395. Parece um passo intermediário de ln 2 colado na sequência. O resultado final (-15$\cdot$0,0693 = -1,0395) está correto.%
\item%
\textit{(statement)} Convenção de velocidade: se v = -$\frac12$$\cdot$d{[}H2O2{]}/dt = k{[}H2O2{]}, a constante efetiva seria 2k e a concentração em 5 min seria 0,25 mol/L. A resolução e o gabarito usam -d{[}H2O2{]}/dt = k{[}H2O2{]}. Essa leitura é apoiada pelo valor tabelado e\^{}(-0,3465), que é justamente e\^{}(-kt) com t = 5 min. A resolução diz que 'a própria prova indica essa convenção', o que é um pouco forte: a prova não diz isso explicitamente, apenas sugere por meio da tabela.%
\item%
\textit{(solution)} Para o item a, a resolução adota hipóteses implícitas razoáveis, todas declaradas: consumo total no 'término' da reação, O2 como gás ideal e seco (sem descontar a pressão de vapor da água) e solubilidade do O2 desprezada. O gabarito adota as mesmas hipóteses sem mencioná{-}las.%
\end{itemize}%
\section*{Questão 3}%
\href{https://poliedroresolve.sistemapoliedro.com.br/vestibulares/ita/2026/ita-2-fase-dia-2-29-10-2025-2-fase-dia-2/dissertativa-3-quimica-2-fase-dia-2-ita-2026}{Resolução do gabarito}\par\smallskip%
\begin{longtable}{@{}p{0.06\textwidth} p{0.33\textwidth} p{0.33\textwidth} p{0.1\textwidth} p{0.08\textwidth}@{}}\toprule Item & Resolução & Gabarito & Veredito & Certo \\ \midrule
a & t $\approx$ 5,5$\times$10$^2$ s (552,8 s com $\surd$5 = 2,236; 552 s com $\surd$5 = 2,24), com V\_p = $\surd$0,2 L $\approx$ 0,4472 L & $\Delta$t = 1240/2,24 s $\approx$ 553,6 s (V' = 1/2,24 L) & equivalente & ambos \\
b & {[}A?{]} = $\surd$5(1+$\surd$3)$\cdot$10$^{-}$$^3$ $\approx$ 6,11$\cdot$10$^{-}$$^3$ mol/L; {[}B$^{-}${]} = $\surd$5($\surd$3-1)$\cdot$10$^{-}$$^3$ $\approx$ 1,64$\cdot$10$^{-}$$^3$ mol/L (com AB(s) formado $\approx$ 0,634$\cdot$10$^{-}$$^3$ mol) & {[}A?{]} = 6,135$\cdot$10$^{-}$$^3$ mol/L; {[}B$^{-}${]} = 1,655$\cdot$10$^{-}$$^3$ mol/L (x = 2,825$\cdot$10$^{-}$$^3$ mol/L) & equivalente & resolução \\
\bottomrule\end{longtable}%
\begin{itemize}%
\item%
\textbf{a)} Os dois usam o mesmo modelo: o volume cai 1 mL$\cdot$s$^{-}$$^{1}$ (540 cal$\cdot$s$^{-}$$^{1}$ ÷ 540 cal$\cdot$g$^{-}$$^{1}$ $\times$ 1 g$\cdot$mL$^{-}$$^{1}$) e a precipitação começa quando 2$\cdot$10$^{-}$$^{6}$/V$^2$ = 10$^{-}$$^{5}$, ou seja, V\_p = 1/$\surd$5 L $\approx$ 0,4472 L. O valor exato é t = 1000(1 - 1/$\surd$5) $\approx$ 552,79 s. O gabarito chega a 553,6 s porque arredonda $\surd$5 para 2,24 e calcula 1/2,24 = 0,4464 L. A resolução escreve 2,24/5 = 0,448 L e obtém 552 s. A diferença entre 552,8 s e 553,6 s vem só do arredondamento de $\surd$5, e as duas respostas dão 5,5$\times$10$^2$ s.%
\item%
\textbf{b)} O método é o mesmo. As concentrações dobram (8,96$\cdot$10$^{-}$$^3$ e 4,48$\cdot$10$^{-}$$^3$ no gabarito), há precipitação e o sistema volta à saturação. Na conta do gabarito, o termo constante (8,96$\cdot$10$^{-}$$^3$)(4,48$\cdot$10$^{-}$$^3$) - 10$^{-}$$^{5}$ vale 3,014$\cdot$10$^{-}$$^{5}$ e foi arredondado para 3$\cdot$10$^{-}$$^{5}$. Com isso x passa de 2,845$\cdot$10$^{-}$$^3$ para 2,825$\cdot$10$^{-}$$^3$, e as concentrações finais saem 6,135$\cdot$10$^{-}$$^3$ e 1,655$\cdot$10$^{-}$$^3$, cujo produto é 1,014$\cdot$10$^{-}$$^{5}$ $\neq$ K\_PS. Sem esse arredondamento, pela mesma via do gabarito, obtêm{-}se 6,115$\cdot$10$^{-}$$^3$ e 1,635$\cdot$10$^{-}$$^3$. Os valores exatos são $\surd$5(1$\pm$$\surd$3)$\cdot$10$^{-}$$^3$, ou seja, 6,109$\cdot$10$^{-}$$^3$ e 1,637$\cdot$10$^{-}$$^3$, com produto exatamente 10$^{-}$$^{5}$. Esses são os valores da resolução. O desvio do gabarito é de cerca de 0,4\% em {[}A?{]} e 1,1\% em {[}B$^{-}${]}. É erro numérico de arredondamento, não conceitual, mas faz {[}B$^{-}${]} com 2 algarismos significativos virar 1,7$\cdot$10$^{-}$$^3$ em vez do correto 1,6$\cdot$10$^{-}$$^3$.%
\end{itemize}%
\paragraph{Pontos de dúvida}%
\begin{itemize}%
\item%
\textit{(key)} No item b, o gabarito Poliedro arredonda o termo independente da equação do 2º grau de 3,014$\cdot$10$^{-}$$^{5}$ para 3$\cdot$10$^{-}$$^{5}$. Isso dá x = 2,825$\cdot$10$^{-}$$^3$ em vez de 2,845$\cdot$10$^{-}$$^3$ e concentrações finais de 6,135$\cdot$10$^{-}$$^3$ e 1,655$\cdot$10$^{-}$$^3$ mol/L, cujo produto (1,014$\cdot$10$^{-}$$^{5}$) não fecha o K\_PS. Os valores corretos são $\approx$ 6,11$\cdot$10$^{-}$$^3$ e 1,64$\cdot$10$^{-}$$^3$ mol/L. A diferença é pequena e não deve mudar a nota, mas {[}B$^{-}${]} com 2 algarismos significativos sai 1,7$\cdot$10$^{-}$$^3$ no gabarito e 1,6$\cdot$10$^{-}$$^3$ no valor correto.%
\item%
\textit{(key)} No item a, o gabarito arredonda $\surd$5 = 2,24 antes de calcular V' = 1/2,24 L e obtém 553,6 s. O valor exato é 552,8 s. Com o mesmo dado ($\surd$5 = 2,24), a resolução usa V\_p = $\surd$5/5 e obtém 552 s. As três versões são aceitáveis e a diferença é só de arredondamento.%
\item%
\textit{(statement)} A consideração ii (ebulição nas condições padrão, 1 bar) contradiz levemente o '100 $^\circ$C' do enunciado (a 1 bar a água ferve a $\approx$ 99,6 $^\circ$C). Gabarito e resolução tratam isso como irrelevante, e ninguém usa o calor sensível. Também fica implícito que o volume do sólido AB formado é desprezível e que não há supersaturação.%
\end{itemize}%
\section*{Questão 4}%
\href{https://poliedroresolve.sistemapoliedro.com.br/vestibulares/ita/2026/ita-2-fase-dia-2-29-10-2025-2-fase-dia-2/dissertativa-4-quimica-2-fase-dia-2-ita-2026}{Resolução do gabarito}\par\smallskip%
\begin{longtable}{@{}p{0.06\textwidth} p{0.33\textwidth} p{0.33\textwidth} p{0.1\textwidth} p{0.08\textwidth}@{}}\toprule Item & Resolução & Gabarito & Veredito & Certo \\ \midrule
a & 5,58 g (cerca de 5,6 g), calculado direto com o resíduo de 28,5\%. Com n = 9, a confirmação cruzada dá 5,56 g. & 5,57 g (usando n = 9, M = 465 g/mol, M média do BGN = 66,8 g/mol e n(Nb2O5) = 1/167 mol) & equivalente & ambos \\
b & n = 9 (M = 465,18 g/mol; resíduo teórico de 28,57\%) & n = 9 & igual & ambos \\
\bottomrule\end{longtable}%
\begin{itemize}%
\item%
\textbf{a)} Refiz a conta. As massas molares são SiO2 = 60,09, CaO = 56,08 e Nb2O5 = 265,82 g/mol. Em 100 mol de óxidos a massa é 6687,56 g, então a fração mássica de Nb2O5 é 0,15899 e 10 g de BGN contêm 1,5899 g de Nb2O5. Pelo caminho da resolução, 1,5899/0,285 = 5,579 g. Pelo caminho do gabarito, com n = 9 (M = 465,18 g/mol), dá 2$\cdot$(1,5899/265,82)$\cdot$465,18 = 5,565 g. Com os arredondamentos do gabarito (66,8 g/mol e 1/167 mol) dá 5,569, ou seja, 5,57 g. A diferença de 0,2\% vem só da escolha entre usar o 28,5\% medido (o resíduo teórico com n = 9 é 28,57\%) e usar o n inteiro. Os dois métodos estão corretos e ambos dão cerca de 5,6 g. A própria resolução mostra os dois valores e explica a diferença.%
\item%
\textbf{b)} Pela razão Nb2O5/(2$\cdot$M\_sal) = 0,285, M\_sal = 265,82/0,57 = 466,35 g/mol. Com M\_anidro = 303,00 g/mol, n = (466,35 - 303,00)/18,02 = 9,06, que arredonda para 9. Os inteiros vizinhos dão resíduos de 29,72\% (n = 8) e 27,51\% (n = 10), longe do dado. Observação: a equação do gabarito, escrita como 0,285 = 266/(303 + 18n), está sem o fator 2 no denominador. Lida ao pé da letra, daria n $\approx$ 35. Mas o texto do gabarito diz que 2 mol de oxalato formam 1 mol de Nb2O5, então é erro de digitação e a resposta n = 9 está certa.%
\end{itemize}%
\paragraph{Pontos de dúvida}%
\begin{itemize}%
\item%
\textit{(key)} No item b, o gabarito escreve 0,285 = 266/(303 + 18n), sem o fator 2 no denominador. O correto é 0,285 = 266/{[}2(303 + 18n){]}. Lida literalmente, a equação dá n $\approx$ 35. O texto do gabarito ('2 mols de oxalato $\to$ 1 mol de Nb2O5') e a resposta n = 9 mostram que é erro de digitação, mas a equação como está escrita não leva ao resultado.%
\item%
\textit{(statement)} O item a) vem antes do b), mas o gabarito usa o n = 9 do item b) para calcular a massa (5,57 g). A resolução usa diretamente o 28,5\% medido (5,58 g). Como o enunciado não diz qual base usar, os dois caminhos são válidos e a diferença (cerca de 0,2\%) é só de arredondamento. Vale confirmar com o professor se a banca aceita os dois.%
\item%
\textit{(solution)} A resolução usa hipóteses que o gabarito também usa sem dizer: a porcentagem da termogravimetria é sobre a massa do sal hidratado, a composição molar se refere às unidades de óxido e todo o Nb passa para o BGN. São todas razoáveis. A observação de que 5,58 g é um 'limite inferior' vai além do que foi pedido, mas não está errada.%
\item%
\textit{(solution)} A equação ilustrativa de decomposição na resolução (produtos N2, CO2 e H2O) é diferente da do gabarito, que cita NH3, H2O, CO2 e CO. Conferi o balanceamento da resolução e ele está correto. Isso não muda o cálculo, que só depende de 2 mol de sal formarem 1 mol de Nb2O5.%
\end{itemize}%
\section*{Questão 5}%
\href{https://poliedroresolve.sistemapoliedro.com.br/vestibulares/ita/2026/ita-2-fase-dia-2-29-10-2025-2-fase-dia-2/dissertativa-5-quimica-2-fase-dia-2-ita-2026}{Resolução do gabarito}\par\smallskip%
\begin{longtable}{@{}p{0.06\textwidth} p{0.33\textwidth} p{0.33\textwidth} p{0.1\textwidth} p{0.08\textwidth}@{}}\toprule Item & Resolução & Gabarito & Veredito & Certo \\ \midrule
a & 3 arranjos, um para cada átomo central: N--C--S (C central), C--N--S (N central) e C--S--N (S central). As formas de ressonância não contam como arranjos novos. & 3 arranjos: {[}CNS{]}$^{-}$, {[}NCS{]}$^{-}$ e {[}CSN{]}$^{-}$ ({[}XGQ{]}$^{-}$, {[}GXQ{]}$^{-}$ e {[}XQG{]}$^{-}$) & igual & ambos \\
b & BF3 (trigonal planar) e NH3 (pirâmide trigonal) & Trigonal planar: BF3, BH3 ou SO3 (ZY3, ZA3, QR3). Pirâmide trigonal: NH3 ou NF3 (GA3, GY3) & igual & ambos \\
c & Nitração do C6H5F por NO2+ (SEAr). Produtos principais: 1{-}fluoro{-}4{-}nitrobenzeno (para, majoritário) e 1{-}fluoro{-}2{-}nitrobenzeno (orto). O meta é desprezível e não há dinitração a 25 $^\circ$C. As estruturas estão desenhadas. & Nitração do C6H5F. O F orienta para orto/para. Produtos: 1{-}fluoro{-}2{-}nitrobenzeno (orto) e 1{-}fluoro{-}4{-}nitrobenzeno (para) & igual & ambos \\
d & H2S < CH3OH < H2O, ou seja, (ii) < (iii) < (i) & H2S < CH3OH < H2O, ou seja, (ii) < (iii) < (i) & igual & ambos \\
\bottomrule\end{longtable}%
\begin{itemize}%
\item%
\textbf{a)} O número e os três arranjos são os mesmos nos dois. Refiz as contas da resolução: são 4+5+6+1 = 16 elétrons de valência. As cargas formais das estruturas citadas estão certas. Com C central, {[}N=C=S{]}$^{-}$ dá $\Sigma$FC = 1, com o -1 no N. Com N central, {[}C$^{-}$$\equiv$N?-S$^{-}${]} dá $\Sigma$ = 3. Com S central, {[}C=S=N{]}$^{-}$ fica C -2, S +2 e N -1, $\Sigma$ = 5. Todas somam -1 e completam o octeto. A conclusão de que o arranjo com C central (tiocianato) é o mais estável está correta. Essa análise vai além do que o item pede, mas não muda a contagem.%
\item%
\textbf{b)} O gabarito aceita várias respostas, e as escolhidas pela resolução (BF3 e NH3) estão entre elas. A justificativa pelo VSEPR está correta: BF3 é AX3 com ângulos de 120$^\circ$, e NH3 é AX3E com ângulo de cerca de 107$^\circ$.%
\item%
\textbf{c)} Os mesmos dois produtos aparecem nos dois. Nos desenhos da resolução, o NO2 está no C4 (para) e no C2 (orto) em relação ao F no C1, como deve ser. A equação HNO3 + 2 H2SO4 $\to$ NO2+ + H3O+ + 2 HSO4- está balanceada (H 5=5, O 11=11, S 2=2, carga 0=0). A regiosseletividade está bem justificada: o F é -I e +M, o para é favorecido e o meta é desprezível. As porcentagens (para $\approx$ 85--90 \%, orto $\approx$ 10--15 \%) batem, de modo geral, com os dados experimentais conhecidos para a nitração do fluorobenzeno.%
\item%
\textbf{d)} A ordem é a mesma, e a justificativa também: as ligações de hidrogênio pesam mais que a massa molar. As massas molares estão certas (18,02, 34,08 e 32,05 g/mol), assim como as temperaturas de ebulição de referência (cerca de -60 $^\circ$C, 65 $^\circ$C e 100 $^\circ$C) e as eletronegatividades de Pauling (O 3,44, S 2,58, H 2,20).%
\end{itemize}%
\paragraph{Pontos de dúvida}%
\begin{itemize}%
\item%
\textit{(statement)} Item a: 'arranjo atômico' tem mais de uma leitura possível. A resolução e o gabarito contam só cadeias abertas e chegam a 3. A resolução lembra que, se o anel triangular C--N--S também contasse, o total seria 4. Vale confirmar com um professor se a banca aceitaria outra interpretação, como contar formas de ressonância ou o anel.%
\item%
\textit{(solution)} Item c: as porcentagens (para $\approx$ 85--90 \%, orto $\approx$ 10--15 \%, meta $\leq$ 1 \%) não vêm do enunciado. São valores aproximados da literatura, que variam um pouco conforme a fonte. Não afetam a resposta, mas não devem ser apresentados como valores exatos.%
\item%
\textit{(solution)} Item c: a afirmação de que 'orto e para recebem a mesma estabilização por ressonância' simplifica o argumento qualitativo. Não está errada no nível do vestibular e não muda a conclusão.%
\end{itemize}%
\section*{Questão 6}%
\href{https://poliedroresolve.sistemapoliedro.com.br/vestibulares/ita/2026/ita-2-fase-dia-2-29-10-2025-2-fase-dia-2/dissertativa-6-quimica-2-fase-dia-2-ita-2026}{Resolução do gabarito}\par\smallskip%
\begin{longtable}{@{}p{0.06\textwidth} p{0.33\textwidth} p{0.33\textwidth} p{0.1\textwidth} p{0.08\textwidth}@{}}\toprule Item & Resolução & Gabarito & Veredito & Certo \\ \midrule
a & Não favorece os produtos. log K1 = pKb(propinido) - pKb(H2O) = -11,00 - 15,70 = -26,70, então K1 $\approx$ 2$\times$10\^{}{-}27. Valores usados: pKa(H3O+) = -1,70 e pKb(CH3CC-) = -11,00. & Não favorece os produtos (deslocado para a esquerda). Keq = 10\^{}(-1,70-25) < 1, com pKa(H3O+) = -1,70 e pKb(CH3CC-) = -11. & igual & ambos \\
b & Não favorece os produtos. log K2 = -11,00 - 4,75 = -15,75, então K2 $\approx$ 1,8$\times$10\^{}{-}16. pKa(NH4+) = 9,25. & Não favorece os produtos (deslocado para a esquerda). pKa(NH4+) = 14 - 4,75 = 9,25 e Keq = 10\^{}(9,25-25) < 1. & igual & ambos \\
c & Favorece os produtos, de forma praticamente completa. log K3 = -11,00 - (-22,00) = +11,00, então K3 = 10\^{}11. pKa(NH3) = 36,00. A resolução ressalva que a reação só ocorre em solvente não aquoso. & Favorece os produtos (deslocado para a direita). pKa(NH3) = 14 - (-22) = 36 e Keq = 10\^{}(36-25) = 10\^{}11 > 1. & igual & ambos \\
\bottomrule\end{longtable}%
\begin{itemize}%
\item%
\textbf{a)} A conclusão e os pK derivados coincidem. Conferi a conta: 10\^{}{-}26,70 $\approx$ 1,995$\times$10\^{}{-}27. A conferência que a resolução faz pela via dos Ka, pKa(H3O+) - pKa(propino) = -26,70, bate com a expressão do gabarito.%
\item%
\textbf{b)} Tudo coincide: 10\^{}{-}15,75 $\approx$ 1,78$\times$10\^{}{-}16 e 9,25 - 25 = -15,75.%
\item%
\textbf{c)} Mesma conclusão e mesmo K = 10\^{}11. A ressalva da resolução sobre o solvente é um acréscimo correto: a reação NH2- + H2O $\to$ NH3 + OH- é muito favorável, por isso acetiletos se preparam com NaNH2 em NH3 líquida. Ela não contradiz o gabarito.%
\end{itemize}%
\paragraph{Pontos de dúvida}%
\begin{itemize}%
\item%
\textit{(statement)} O dado pKb(H2O) = 15,70 parte da convenção Kw/{[}H2O{]} = Kw/55,5. Ela é termodinamicamente inconsistente com pKw = 14, e o próprio Poliedro critica o valor (J. Chem. Educ. 2017, 94, 690). Com pKb(H2O) = 14 o resultado seria pKa(H3O+) = 0 e log K1 = -25, e a conclusão do item a) seria a mesma. Não muda nenhuma resposta.%
\item%
\textit{(solution)} A resolução acrescenta, por hipótese própria, pKa(H2O) = 15,70 e pKb(OH-) = -1,70. Esses valores não são dados no enunciado nem estão no gabarito. Pela convenção pKa(H2O) = pKw = 14, seria pKb(OH-) = 0. Como não entram nas respostas de a), b) e c), não afetam a nota. Servem apenas para a ressalva de solvente do item c), cuja conclusão vale nas duas convenções.%
\item%
\textit{(solution)} A forma genérica HA + B ? A- + BH+ do Passo 3 tem cargas inconsistentes quando B = NH2- (reação 3). A relação K = Kb(B)/Kb(A-) continua válida, então é só imprecisão de notação.%
\end{itemize}%
\section*{Questão 7}%
\href{https://poliedroresolve.sistemapoliedro.com.br/vestibulares/ita/2026/ita-2-fase-dia-2-29-10-2025-2-fase-dia-2/dissertativa-7-quimica-2-fase-dia-2-ita-2026}{Resolução do gabarito}\par\smallskip%
\begin{longtable}{@{}p{0.06\textwidth} p{0.33\textwidth} p{0.33\textwidth} p{0.1\textwidth} p{0.08\textwidth}@{}}\toprule Item & Resolução & Gabarito & Veredito & Certo \\ \midrule
a & Meio ácido: cátodo 2H? + 2e$^{-}$ $\to$ H$_{2}$; ânodo H$_{2}$O $\to$ $\frac12$O$_{2}$ + 2H? + 2e$^{-}$. Também dá a versão em meio básico: cátodo 2H$_{2}$O + 2e$^{-}$ $\to$ H$_{2}$ + 2OH$^{-}$; ânodo 2OH$^{-}$ $\to$ $\frac12$O$_{2}$ + H$_{2}$O + 2e$^{-}$. & Cátodo: 2H$_{2}$O + 2e$^{-}$ $\to$ H$_{2}$ + 2OH$^{-}$; ânodo: 2H$_{2}$O $\to$ O$_{2}$ + 4H? + 4e$^{-}$ & equivalente & ambos \\
b & 5,32$\times$10$^{5}$ J (532 kJ), que é a energia de 2 mol de fótons de 450 nm & 5,304$\times$10$^{5}$ J = 530,4 kJ (2 mol de fótons) & equivalente & ambos \\
c & 1,35$\times$10$^{-}$$^2$ mol de H$_{2}$ & $\approx$1,36$\times$10$^{-}$$^2$ mol de H$_{2}$ & equivalente & ambos \\
d & 1,49$\times$10$^{9}$ J (2797 mol de H$_{2}$ $\times$ 2 mol de fótons de 450 nm) & $\approx$1,5 GJ (1 498 380 kJ, com 2825 mol de H$_{2}$) & equivalente & ambos \\
e & $\lambda$max $\approx$ 837 nm (143 kJ por mol de fótons; 2 fótons por H$_{2}$) & $\lambda$ $\approx$ 835 nm & equivalente & ambos \\
\bottomrule\end{longtable}%
\begin{itemize}%
\item%
\textbf{a)} As duas respostas trazem semirreações corretas de redução da água (ou do H?) a H$_{2}$ e de oxidação da água a O$_{2}$. A da resolução está toda em meio ácido e confere a soma com a equação global. O gabarito mistura os meios: cátodo em meio básico e ânodo em meio ácido. Mesmo assim a soma dá a equação global, porque 4H? + 4OH$^{-}$ formam 4H$_{2}$O. A resolução ainda apresenta a forma em meio básico, que coincide com o cátodo do gabarito. Mudam só o meio e os coeficientes.%
\item%
\textbf{b)} O método é o mesmo: 2 mol de e$^{-}$, portanto 2 mol de fótons, e E = 2$\cdot$N\_A$\cdot$hc/$\lambda$. A conta refeita dá 2$\cdot$6,02$\times$10$^2$$^3$$\cdot$6,63$\times$10$^{-}$$^3$$^{4}$$\cdot$3,0$\times$10$^{8}$/4,50$\times$10$^{-}$$^{7}$ = 5,32$\times$10$^{5}$ J. O gabarito chega a 5,304$\times$10$^{5}$ J porque arredonda N\_A para 6$\times$10$^2$$^3$, em vez dos 6,02$\times$10$^2$$^3$ dados pela prova. A diferença, de 0,3\%, vem só desse arredondamento. A resolução usa a constante dada e fica mais fiel aos dados.%
\item%
\textbf{c)} Energia incidente: 10$\cdot$0,2$\cdot$3600 = 7200 J nas duas soluções. Dividindo pelos 5,32$\times$10$^{5}$ J por mol de H$_{2}$, obtém{-}se 1,353$\times$10$^{-}$$^2$ mol. Com o valor do gabarito, 5,304$\times$10$^{5}$ J, obtém{-}se 1,357$\times$10$^{-}$$^2$ mol. A diferença vem apenas do N\_A arredondado no gabarito.%
\item%
\textbf{d)} As duas soluções interpretam a pergunta da mesma forma: a energia luminosa é a dos fótons de 450 nm absorvidos pelo semicondutor, multiplicada pelos mols de H$_{2}$. Os mols mudam: a resolução usa M(H$_{2}$) = 2,02 g/mol, dado da prova, e obtém 5650/2,02 = 2797 mol. O gabarito usa 2 g/mol e obtém 2825 mol, além de N\_A = 6$\times$10$^2$$^3$. A conta refeita dá 2797$\cdot$5,32$\times$10$^{5}$ = 1,488$\times$10$^{9}$ J, contra 1,498$\times$10$^{9}$ J do gabarito. Arredondados, os dois valores dão 1,5 GJ.%
\item%
\textbf{e)} O raciocínio é o mesmo nas duas soluções. A combustão libera 286 kJ por mol de H$_{2}$, e cada mol de H$_{2}$ exige 2 mol de fótons. Portanto 2$\cdot$N\_A$\cdot$hc/$\lambda$ = 286$\times$10$^3$ J. Com N\_A = 6,02$\times$10$^2$$^3$, $\lambda$ = 837,3 nm, que é o valor da resolução. Com N\_A = 6$\times$10$^2$$^3$, $\lambda$ = 834,5 nm $\approx$ 835 nm, que é o valor do gabarito. A diferença, de 0,3\%, vem só da constante arredondada pelo gabarito.%
\end{itemize}%
\paragraph{Pontos de dúvida}%
\begin{itemize}%
\item%
\textit{(key)} O gabarito do Poliedro usa N\_A = 6$\times$10$^2$$^3$ mol$^{-}$$^{1}$ e M(H$_{2}$) = 2 g/mol. A prova fornece N\_A = 6,02$\times$10$^2$$^3$ mol$^{-}$$^{1}$ e M(H) = 1,01 g/mol, ou seja, M(H$_{2}$) = 2,02 g/mol. Por isso os valores do gabarito diferem um pouco: 530,4 kJ contra 532 kJ, 2825 mol contra 2797 mol, 835 nm contra 837 nm. A resolução usa as constantes dadas. As diferenças ficam entre 0,3\% e 1\% e não mudam o método.%
\item%
\textit{(key)} No item (a), o gabarito escreve o cátodo em meio básico (2H$_{2}$O + 2e$^{-}$ $\to$ H$_{2}$ + 2OH$^{-}$) e o ânodo em meio ácido (2H$_{2}$O $\to$ O$_{2}$ + 4H? + 4e$^{-}$). Isso é inconsistente quanto ao meio, embora a soma ainda resulte na equação global. Cada semirreação está correta isoladamente, e a resolução é consistente: dá tudo em meio ácido e também a versão em meio básico.%
\item%
\textit{(statement)} No item (d), a expressão "energia mínima necessária {[}...{]} utilizando o semicondutor" admite outra leitura. Se "mínima" fosse entendida como a energia termodinâmica, a resposta seria $\Delta$H = n$\cdot$286 kJ $\approx$ 8,0$\times$10$^{8}$ J. O gabarito e a resolução adotam a mesma leitura: a energia dos fótons de 450 nm, $\approx$ 1,5$\times$10$^{9}$ J. A resolução discute e justifica explicitamente a escolha. Vale confirmar com o professor se a banca aceitaria a outra leitura.%
\item%
\textit{(statement)} No item (e), as duas soluções igualam a energia luminosa ao |$\Delta$H| de combustão com formação de H$_{2}$O(?) (286 kJ/mol), e 2 fótons por H$_{2}$, como o enunciado pede. O critério termodinâmico rigoroso seria $\Delta$G$^\circ$, de cerca de 237 kJ/mol, que não é fornecido. A resolução comenta esse ponto no item (b).%
\end{itemize}%
\section*{Questão 8}%
\href{https://poliedroresolve.sistemapoliedro.com.br/vestibulares/ita/2026/ita-2-fase-dia-2-29-10-2025-2-fase-dia-2/dissertativa-8-quimica-2-fase-dia-2-ita-2026}{Resolução do gabarito}\par\smallskip%
\begin{longtable}{@{}p{0.06\textwidth} p{0.33\textwidth} p{0.33\textwidth} p{0.1\textwidth} p{0.08\textwidth}@{}}\toprule Item & Resolução & Gabarito & Veredito & Certo \\ \midrule
a & CH4(g) + 2 O2(g) $\to$ CO2(g) + 2 H2O(g); $\Delta$Hf{[}H2O(g){]} = -241 kJ/mol; $\Delta$Hc = -800 kJ/mol & CH4(g) + 2 O2(g) $\to$ CO2(g) + 2 H2O(g), $\Delta$H = -800 kJ/mol de CH4 ($\Delta$Hf H2O(g) = -241). Para 12 MJ: 15 CH4 + 30 O2 $\to$ 15 CO2 + 30 H2O(g), $\Delta$H = -12 000 kJ & equivalente & ambos \\
b & m(CH4) = 240,75 g ($\approx$ 241 g); m(O2) = 960 g & m(CH4) = 240 g; m(O2) = 960 g & equivalente & ambos \\
c & $\approx$ 0,644 g de CH4(g) (0,64 g com 2 algarismos significativos) & 0,64 g (Q útil = 240$\cdot$40 J = 9,6 kJ; Q de combustão na pré{-}câmara = 32 kJ) & equivalente & ambos \\
d & 200 g de CH4(?) (massa mínima) & 200 g (Q = 0,25\%$\cdot$12 000 kJ = 30 kJ = m$\cdot$(5/3)$\cdot$90 K; mínima por usar a temperatura máxima de 188 K) & igual & ambos \\
\bottomrule\end{longtable}%
\begin{itemize}%
\item%
\textbf{a)} A equação e o $\Delta$H por mol (-800 kJ/mol) são os mesmos. Confirmei: {[}-393 + 2(-286+45){]} - (-75) = -800. O gabarito também escreve a equação multiplicada por 15, com $\Delta$H = -12 MJ. A resolução não faz isso no item a, mas obtém os 15 mol logo no item b. O conteúdo é o mesmo, só a apresentação muda.%
\item%
\textbf{b)} A diferença vem só da massa molar usada. O gabarito arredonda M(CH4) para 16. A resolução usa os dados da prova, 12,01 + 4$\cdot$1,01 = 16,05 g/mol, e chega a 15$\cdot$16,05 = 240,75 g, valor mais fiel aos dados. Para o O2, os dois dão 30$\cdot$32,00 = 960 g.%
\item%
\textbf{c)} O raciocínio é o mesmo nos dois: 40 J/g vezes a massa pressurizada da câmara principal dá o Q útil, que é 30\% do calor da pré{-}câmara, a 800 kJ/mol. Gabarito: 9600 J / 0,3 = 32 kJ $\to$ 0,04 mol $\to$ 0,64 g. Resolução: 9630 J / 0,3 = 32,1 kJ $\to$ 0,040125 mol $\to$ 0,644 g. A diferença vem só do arredondamento de M (16 contra 16,05). Conferi também a leitura alternativa da resolução, em que a pré{-}câmara está incluída nos 15 mol: dá 0,642 g. Todos os valores arredondam para 0,64 g.%
\item%
\textbf{d)} Os valores e o argumento de mínimo coincidem: Q = 0,0025$\cdot$1,2$\times$10\^{}7 J = 3,0$\times$10\^{}4 J e m = 3,0$\times$10\^{}4 / {[}(5/3)$\cdot$90{]} = 200 g. Os dois justificam o mínimo pela $\Delta$T máxima de 90 K.%
\end{itemize}%
\paragraph{Pontos de dúvida}%
\begin{itemize}%
\item%
\textit{(solution)} Item a: o pedido é 'calcule a entalpia da combustão, conforme o requisito I'. O gabarito entende isso como $\Delta$H = -12 MJ para a equação com 15 CH4 + 30 O2. A resolução dá só -800 kJ/mol no item a e chega aos 15 mol apenas no item b. Uma banca exigente poderia querer ver a equação multiplicada por 15 (ou $\Delta$H = -12 MJ) já no item a.%
\item%
\textit{(key)} Item b: o gabarito usa M(CH4) = 16 g/mol e obtém 240 g. Com as massas molares dadas na prova (C 12,01; H 1,01), M = 16,05 g/mol e a massa é 240,75 g. É só arredondamento, mas o valor da resolução é o mais fiel aos dados fornecidos.%
\item%
\textit{(statement)} Item c: o enunciado deixa duas coisas em aberto. A primeira é se a 'pequena porcentagem do combustível' da pré{-}câmara faz parte dos 15 mol que liberam 12 MJ. A segunda é qual massa a turbobomba pressuriza: só o metano da câmara principal ou também o de outras linhas. O gabarito e a resolução adotam 'só a câmara principal, 240 g'. A leitura alternativa (pré{-}câmara dentro dos 15 mol) dá 0,642 g. As duas leituras arredondam para 0,64 g, então a nota não muda.%
\item%
\textit{(solution)} Item d: o 'fechamento' da resolução diz que 200 g passam pelo resfriamento e 40,75 g vão direto à injeção. Pelo esquema, todo o CH4(?) que sai da turbobomba parece seguir pela linha de refrigeração. A frase é uma interpretação extra, que não estava no pedido e não afeta o valor mínimo de 200 g. A hipótese de que o metano 'permanece líquido' até 188 K também é só uma idealização do enunciado.%
\end{itemize}%
\section*{Questão 9}%
\href{https://poliedroresolve.sistemapoliedro.com.br/vestibulares/ita/2026/ita-2-fase-dia-2-29-10-2025-2-fase-dia-2/dissertativa-9-quimica-2-fase-dia-2-ita-2026}{Resolução do gabarito}\par\smallskip%
\begin{longtable}{@{}p{0.06\textwidth} p{0.33\textwidth} p{0.33\textwidth} p{0.1\textwidth} p{0.08\textwidth}@{}}\toprule Item & Resolução & Gabarito & Veredito & Certo \\ \midrule
a & Principais: Cu$^2$?, K?, H?, SO$_{4}$$^2$$^{-}$, HSO$_{4}$$^{-}$, Cr$_{2}$O$_{7}$$^2$$^{-}$ e HCrO$_{4}$$^{-}$; traços: CrO$_{4}$$^2$$^{-}$ e OH$^{-}$. & Principais: H? (H$_{3}$O?), SO$_{4}$$^2$$^{-}$, Cu$^2$?, K?, Cr$_{2}$O$_{7}$$^2$$^{-}$; em quantidades muito pequenas: OH$^{-}$, HSO$_{4}$$^{-}$, CrO$_{4}$$^2$$^{-}$. & equivalente & ambos \\
b & Migração iônica: o Cu$^2$? (azul) migra para o cátodo (Pt$_{1}$, braço esquerdo) e o Cr$_{2}$O$_{7}$$^2$$^{-}$ (laranja) para o ânodo (Pt$_{2}$, braço direito). Os íons coloridos não se descarregam porque os eletrodos estão na zona de H$_{2}$SO$_{4}$. Após um tempo finito, a migração é parcial, o que forma as faixas; o fundo continua marrom (mistura original). & Migração iônica: Pt$_{1}$ (negativo) atrai os cátions K?, H$_{3}$O? e Cu$^2$? (azul), e a solução perto de Pt$_{1}$ fica azul. Pt$_{2}$ (positivo) atrai os ânions SO$_{4}$$^2$$^{-}$ e Cr$_{2}$O$_{7}$$^2$$^{-}$ (alaranjado), e a solução perto de Pt$_{2}$ fica alaranjada. & igual & ambos \\
c & Pt$_{1}$ = cátodo: 2H? + 2e$^{-}$ $\to$ H$_{2}$. Pt$_{2}$ = ânodo: 2H$_{2}$O $\to$ O$_{2}$ + 4H? + 4e$^{-}$. Global 2H$_{2}$O $\to$ 2H$_{2}$ + O$_{2}$, $\Delta$G$^\circ$ $\approx$ +4,75$\times$10$^{5}$ J/mol, tensão mínima de 1,23 V. & Pt$_{1}$ é o catodo (redução): 2H?(aq) + 2e$^{-}$ $\to$ H$_{2}$(g). Pt$_{2}$ é o anodo (oxidação): 2H$_{2}$O(l) $\to$ O$_{2}$(g) + 4H?(aq) + 4e$^{-}$. As placas ficam em contato com a solução de H$_{2}$SO$_{4}$. & igual & ambos \\
d & P$_{1}$ é negativo e P$_{2}$ é positivo. Corrente convencional: P$_{2}$ $\to$ Pt$_{2}$ $\to$ solução (da direita para a esquerda) $\to$ Pt$_{1}$ $\to$ P$_{1}$ e, dentro da fonte, de P$_{1}$ para P$_{2}$. Elétrons nos fios: de P$_{1}$ para Pt$_{1}$ e de Pt$_{2}$ para P$_{2}$. & P$_{1}$ é o polo negativo e P$_{2}$ o positivo. Elétrons (sentido real) de P$_{1}$ para P$_{2}$; corrente convencional de P$_{2}$ para P$_{1}$ pelo circuito externo. & igual & ambos \\
\bottomrule\end{longtable}%
\begin{itemize}%
\item%
\textbf{a)} A resolução lista todas as espécies do gabarito e acrescenta o HCrO$_{4}$$^{-}$. A única diferença é de classificação: o gabarito põe o HSO$_{4}$$^{-}$ entre as espécies em quantidade muito pequena, e a resolução o põe entre as principais. Pela química, a resolução está mais próxima do correto. Com Ka2 $\approx$ 1,2$\times$10$^{-}$$^2$, o HSO$_{4}$$^{-}$ é cerca de 90\% do sulfato total em H$_{2}$SO$_{4}$ 0,1 mol/L e cerca de 55\% a 0,01 mol/L. Ele só fica minoritário em soluções muito diluídas, por volta de 10$^{-}$$^3$ mol/L. A conta do HCrO$_{4}$$^{-}$ também confere: com K = 2$\times$10$^{-}$$^2$ e dicromato a 10$^{-}$$^2$ mol/L, dá {[}HCrO$_{4}$$^{-}${]} = 1,0$\times$10$^{-}$$^2$ mol/L, ou seja, 50\% do Cr(VI). As duas listas são aceitáveis para o que o item pede (listar as espécies iônicas).%
\item%
\textbf{b)} O mecanismo é o mesmo (migração dos cátions para o polo negativo/cátodo e dos ânions para o polo positivo/ânodo), com a mesma atribuição de cores e lados. A resolução detalha mais: diz que não há descarga dos íons coloridos, que a migração é parcial e que a fronteira depende das mobilidades. Nada disso contradiz o gabarito.%
\item%
\textbf{c)} A atribuição dos eletrodos e as semirreações são idênticas. A conta extra da resolução está certa: $\Delta$G$^\circ$ = -nFE = -4 $\times$ 9,65$\times$10$^{4}$ $\times$ (-1,23) = +4,748$\times$10$^{5}$ J por mol de O$_{2}$. Os potenciais citados também estão corretos: E(H?/H$_{2}$) = -0,059 pH, E(O$_{2}$/H$_{2}$O) = 1,23 - 0,059 pH e E$^\circ$(S$_{2}$O$_{8}$$^2$$^{-}$) = 2,01 V.%
\item%
\textbf{d)} A polaridade e o sentido são os mesmos. Pela Figura (A), o fio de P$_{1}$ (à esquerda) vai a Pt$_{1}$ e o de P$_{2}$ (à direita) vai a Pt$_{2}$, sem cruzamento, o que confirma P$_{1}$ negativo. A resolução descreve o circuito completo, incluindo o trecho dentro da fonte, de forma coerente com o gabarito.%
\end{itemize}%
\paragraph{Pontos de dúvida}%
\begin{itemize}%
\item%
\textit{(key)} No item a, o gabarito classifica o HSO$_{4}$$^{-}$ como espécie em 'quantidades muito pequenas'. Com Ka2(HSO$_{4}$$^{-}$) $\approx$ 1,2$\times$10$^{-}$$^2$, em H$_{2}$SO$_{4}$ diluído de concentração usual (0,01 a 1 mol/L) o HSO$_{4}$$^{-}$ fica entre cerca de 55\% e 99\% do sulfato total. Ele só é minoritário abaixo de cerca de 10$^{-}$$^3$ mol/L. Isso não muda a nota, porque as duas respostas listam o HSO$_{4}$$^{-}$, mas a classificação do gabarito é discutível.%
\item%
\textit{(key)} O gabarito não menciona o HCrO$_{4}$$^{-}$. Em meio ácido, com dicromato a cerca de 10$^{-}$$^2$ mol/L, aproximadamente metade do Cr(VI) está nessa forma (a conta confere: 50\%). Para um item que pede 'todas as espécies iônicas', a lista da resolução é mais completa.%
\item%
\textit{(solution)} A resolução usa constantes que a prova não fornece (Ka2 do HSO$_{4}$$^{-}$, constante de dimerização do cromato, pKa do HCrO$_{4}$$^{-}$, condutividades iônicas) e uma concentração ilustrativa de 10$^{-}$$^2$ mol/L. Ela declara que são apenas ordem de grandeza e que a conclusão não depende delas. Está correto, mas vai além do que a questão exige.%
\item%
\textit{(statement)} A identificação de P$_{1}$/P$_{2}$ com Pt$_{1}$/Pt$_{2}$ depende de ler na Figura (A) que os fios não se cruzam: P$_{1}$ (esquerda) liga{-}se a Pt$_{1}$ e P$_{2}$ (direita) a Pt$_{2}$. A figura original confirma essa leitura, e tanto a resolução quanto o gabarito a usam.%
\end{itemize}%
\section*{Questão 10}%
\href{https://poliedroresolve.sistemapoliedro.com.br/vestibulares/ita/2026/ita-2-fase-dia-2-29-10-2025-2-fase-dia-2/dissertativa-10-quimica-2-fase-dia-2-ita-2026}{Resolução do gabarito}\par\smallskip%
\begin{longtable}{@{}p{0.06\textwidth} p{0.33\textwidth} p{0.33\textwidth} p{0.1\textwidth} p{0.08\textwidth}@{}}\toprule Item & Resolução & Gabarito & Veredito & Certo \\ \midrule
a & X = 226 & X = 226 & igual & ambos \\
b & Z = 86 (A = Rn{-}226, supondo que J só emite alfa) & Z = 86 & igual & ambos \\
c & 3 partículas alfa & 3 partículas alfa & igual & ambos \\
\bottomrule\end{longtable}%
\begin{itemize}%
\item%
\textbf{a)} Recalculado. Q (232, 89) emite $\beta$$^{-}$ e vira G (232, 90), com N\_G = 142. E é isótono de G e tem Z = 84, então M\_E = 84 + 142 = 226. Como A é isóbaro de E, X = 226.%
\item%
\textbf{b)} Pela conservação de carga, 92 = 2$\cdot$3 + Z, logo Z = 86 (radônio). A hipótese de que J só emite alfa é a leitura literal do enunciado ("J sofre decaimentos de partículas alfa"). O gabarito usa a mesma hipótese sem dizer isso.%
\item%
\textbf{c)} Pela conservação do número de massa, 238 = 4n + 226, logo n = 3. Conferido em Python.%
\end{itemize}%
\paragraph{Pontos de dúvida}%
\begin{itemize}%
\item%
\textit{(statement)} O enunciado diz "massa atômica igual a 238" e "massa 232". A resolução e o gabarito leem esses valores como números de massa inteiros. É a leitura usual e não muda nada nas contas.%
\item%
\textit{(statement)} O valor de Z depende de supor que a cadeia de J só tem emissões alfa. Com k emissões $\beta$$^{-}$ teríamos Z = 86 + k. O texto ("sofre decaimentos de partículas alfa") sustenta a hipótese k = 0. A resolução declara essa hipótese e o gabarito a usa implicitamente.%
\end{itemize}%
\end{document}